% ANCHOR: after_label def:c12-gens
\begin{remark}[Physics notation: symmetry generators and their finite renormalisation]
In the gauge-fixed language $S^{(1)}(x)$ is the source term of the classical transformation $\delta_x$: each field $\varphi$ is multiplied by the variation $\delta_x\varphi$ and paired with its BRST source $\varphi^*$, as in the Zinn-Justin form of the action, so that $\{S^{(1)}(x),\varphi\}=\delta_x\varphi$. For the dilatation $\delta_\Delta\varphi=(2x^m\partial_m+2w_\varphi)\varphi$ with the Weyl weights $w=1$ for $A_m$, $A_I$, $w=\tfrac32$ for $\lambda_\alpha{}^i$, $\bar\lambda_{\dot\alpha i}$ and $w=0$ for $c$ (physics dilatation $D=-\tfrac12\Delta$, key), for a translation $\delta_{P_b}\varphi=b^m\partial_m\varphi$, and for a supercharge $\delta_\xi$ is the supersymmetry transformation with parameters $\xi_\alpha{}^i$, $\bar\xi_{\dot\alpha i}$ of $Q^i_\alpha$, $\bar Q_{\dot\alpha i}$. The one-loop finite counterterm $\hbar F_1^{(1)}(x)$ is the finite local renormalisation of the transformation: a term (local function)$\cdot\varphi^*$ adds that function to $\delta_x\varphi$ at order $g^2N$. The sector parts are the classical transformation restricted to one field type: the six scalars $A_I$, the gauginos, and the gauge field together with the ghost. A coefficient $\theta^x_d$ rescales the classical transformation of one field type by $1+\theta^x_d$ units, and a uniform coefficient $t_x$ rescales the whole generator, $\delta_x\mapsto(1+t_x\,\text{units})\,\delta_x$, with the unit $-\tfrac{g^2N}{32\pi^2}$.
\end{remark}

% ANCHOR: after_label lem:c12-wz
\begin{remark}[Physics notation: anomalies of the regulated Ward identities and their consistency]
$\mathrm{Ob}^{(1)}(x)$ is the one-loop breaking of the Ward identity of the generator $x$ by the regulator, a local operator inserted in the generating functional of the regulated theory, and $\mathrm{Ob}^{(0)}$ the breaking of the BRST identity (block after Theorem \ref{thm:c4-obstruction}). The identities \eqref{eq:c12-wz} are the Wess--Zumino consistency conditions of these breakings: the BRST variation of the breaking of $\delta_x$ is the $\delta_x$-variation of the breaking of BRST, and the breakings of two generators are related through their commutator. When the breakings are BRST variations of local terms, finite counterterms restore the identities; this is the equation $s(\hbar F_1)=\hbar\,\mathrm{Ob}$ read arity by arity. The signs in \eqref{eq:c12-wz} follow from the bracket of Definition \ref{def:c1-bracket}, which differs from the standard antibracket by $(-1)^{|f|}$.
\end{remark}

% ANCHOR: after_label prop:c12-poles
\begin{remark}[Physics notation: the divergences of the symmetry insertions]
With the proper-time cutoff $\varepsilon$ playing the role of $\Lambda_{\rm UV}^{-2}$, the counterterm \eqref{eq:c12-poles} contains one power divergence, the quadratically divergent mass term $\tfrac94\varepsilon^{-1}\int(A,A)$ units of the gauge field and of the six scalars (Theorem \ref{thm:c4-quadratic}), and logarithmic divergences that are field renormalisations, $\log\varepsilon$ times the BRST variation of the field-counting functional $\mathcal N$ (Theorem \ref{thm:c4-log}). Insertions of one or two symmetry transformations carry no power divergence: by dimensional analysis every candidate counterterm would need more fields than its weight allows. Their logarithmic divergences are again field renormalisations of the transformations, up to the closed combination \eqref{eq:c12-closedlog}.
\end{remark}

% ANCHOR: after_label thm:c12-classes
\begin{remark}[Physics notation: the breaking as graphs with one regulator insertion]
The graphs (C) are one-loop graphs with one insertion of the symmetry transformation $\delta_x$ and one regulator insertion, the line $e^{-\varepsilon p^2}-1$ that joins a field to its BRST source (block after Theorem \ref{thm:c4-obstruction}). The regulator insertion lives at proper times below $\varepsilon$, so each such graph is a pure short-distance quantity: powers $\varepsilon^{-m}$ and a finite local polynomial in the momenta, with no $\log\varepsilon$. The finite parts summed over the graphs are the breaking of the Ward identity of $\delta_x$, in the same way in which an anomaly is computed as the finite part of a regulator insertion. The placements of Table \ref{tab:c12-place} list where $\delta_x$ can sit: on the internal ghost--gauge line (for $\Delta$, $P_m$, $M_{mn}$, $K_m$) or gaugino--ghost line (for $Q$, $S$), next to the regulator insertion, or, for the supersymmetry variation $\delta_\xi\lambda\ni[A,A]\xi$, as a cubic vertex.
\end{remark}

% ANCHOR: after_label prop:c12-xi2
\begin{remark}[Physics notation: the supersymmetry Ward identity on two fields]
On two fields the supersymmetry variations are $\delta_\xi A_m=i\bar\lambda_{\dot\alpha i}\bar\sigma^{m\dot\alpha\alpha}\xi_\alpha{}^i$ (the term $T_1$), $\delta_\xi A_{kl}$ from $\xi\lambda$ (the term $T_2$), and $\delta_\xi\lambda$ from $\partial A_m$ and $\partial A_{ij}$ (the terms $T_3$, $T_4$); the two-component forms are in Remark \ref{rem:c12-twocomp}. The breakings $g_1$, $g_3$ are the variations of the gauge and scalar kinetic terms that a rescaling of $\delta_\xi A_m$ or $\delta_\xi A_{ij}$ produces, and $g_2$, $g_4$ involve $\partial^mA_m$ and the ghost, the gauge-fixing sector. $T_0=-\int(\partial^mA_m)\,\xi\lambda^*$ is a modification of $\delta_\xi\lambda$ by the gauge-fixing function; its coefficient is $0$ in the solution, the breaking in $g_2$ being exactly that produced by the finite gauge-fixing counterterm $\tfrac{27}{16}\int(\partial^mA_m)^2$ units. The two free directions rescale the gauge multiplet arrows $A_m\leftrightarrow\lambda$ and the scalar arrows $A_{ij}\leftrightarrow\lambda$ of $\delta_\xi$ separately.
\end{remark}

% ANCHOR: after_label thm:c12-sol
\begin{remark}[Physics notation: the one-loop supersymmetry transformation]
Read through the arrows, $\hbar F_1^{(1)}(\xi)|_2=(\log\tfrac32-\tfrac{79}{12})(T_1+T_2)$ units $+\,t_\xi\,S^{(1)}_2(\xi)$ units multiplies the variations $\delta_\xi A_m$ and $\delta_\xi A_{ij}$ by $1+(\log\tfrac32-\tfrac{79}{12}+t_\xi)$ units and the variation $\delta_\xi\lambda$ by $1+t_\xi$ units, with the unit $-\tfrac{g^2N}{32\pi^2}$: at one loop the regulated theory is supersymmetric with these rescaled transformations. Along the dilatation the particular part $\tfrac{22}3\int(A,A^*)-\tfrac{47}3\int(\lambda^*,\lambda)-\tfrac83\int(c^*,c)$ units is a finite field-strength renormalisation entering the one-loop scale transformation, and the free part $\theta^\Delta_d$ multiplies the classical scale transformation of each field type. Along $K_m$ the solution contains, besides the analogous terms, four second-order symmetries of the free fields, built from the conformal Killing tensors, whose coefficients are fixed at three legs.
\end{remark}

% ANCHOR: after_label thm:c12-fix
\begin{remark}[Physics notation: the one-loop superconformal algebra]
With the rescaled generators $(1+t_x\,\text{units})\delta_x$, the arity-two equations are the one-loop forms of the commutators $[D,P_m]=P_m$, $[D,Q]=\tfrac12Q$, $\{Q^i_\alpha,\bar Q_{\dot\alpha j}\}\propto\delta^i_j\sigma^m_{\alpha\dot\alpha}P_m$, $\{Q,S\}=D+M+R$, $[K,P]=D+M$ and $\{S,\bar S\}\propto K$ (physics dilatation $D=-\tfrac12\Delta$), with the breaking $\varrho$ on the right. Two rescalings are invisible to all of them: the inner one, which multiplies each generator by $1+c\lambda_x$ units, $\lambda_x$ being $-2$ times its scaling dimension, and is the effect of a finite renormalisation of the operators by a multiple of their dimension; and the outer one, which multiplies $Q^i_\alpha$, $\bar S^{\dot\alpha i}$ by $1+a$ units and $\bar Q_{\dot\alpha i}$, $S^\alpha{}_i$ by $1-a$ units, the automorphism of $\mathfrak{psl}(4|4)$ counting upper minus lower $SU(4)$ indices on the odd generators, which is not generated by an element of the algebra.
\end{remark}

% ANCHOR: after_label thm:c12-ar2
\begin{remark}[Physics notation: what the one-loop algebra fixes]
$t_\Delta=-\tfrac13$ states that the finite counterterm of the dilatation contains $-\tfrac13$ units $=\tfrac{g^2N}{96\pi^2}$ times the classical scale transformation of every field type; it is the condition that the one-loop dilatation of an operator at $x$ has no term proportional to $x^m$ times a one-loop translation, the one-loop form of $[D,P_m]=P_m$, and it is reproduced by $[D,Q]=\tfrac12Q$. $t_M=t_\Delta$ and $t_R=t_\Delta+\tfrac13$ fix the rotations relative to the dilatation. The relation $t_{\xi_+}+t_{\xi_-}-t_P=\tfrac{27}4+2\log\tfrac32$ is the one-loop part of $\{Q,\bar Q\}\propto P$: the product of the one-loop normalisations of $Q$ and $\bar Q$ divided by that of $P$, and $t_{\xi_+}+t_{\eta_+}-t_\Delta$ the same for $\{Q,S\}$ against $D$. In a scheme in which the translations are the classical ones, $t_P=-\tfrac13$, the one-loop normalisations of $Q$ and $\bar Q$ multiply to $1+(\tfrac{77}{12}+2\log\tfrac32)$ units. The individual normalisations of $Q$, $\bar Q$, $S$, $\bar S$, $K$ are conventions.
\end{remark}

% ANCHOR: after_label thm:c12-master
\begin{remark}[Physics notation: the one-loop generators on composite operators]
For a composite operator $O$ built from fields $X_1\cdots X_n$, the one-loop correction of the generator $\delta_x$ acting on $O$ is a sum over the pairs of constituent fields of a local two-body density, the finite part of the one-loop graphs through the two fields with one $\delta_x$ insertion and one regulator insertion, plus a one-body term, the finite renormalisation $\hbar F_1^{(1)}(x)$ of $\delta_x$ acting on each field. For the dilatation this is the structure of the one-loop dilatation operator as a sum of densities on neighbouring letters [B03c, (2.12)], here for every generator of $\mathfrak{psl}(4|4)$ and for gauge-variant operators as well. In Chapters 10 and 11 a finite vertex built from the regulator insertion on an external line and the counterterm was added to the density; the graph with the regulator insertion on that external line cancels it exactly (Lemma \ref{lem:c12-bridge}), so the density is the finite part of the graphs through the two fields alone.
\end{remark}

% ANCHOR: after_label thm:c12-D1
\begin{remark}[Physics notation: the anomalous dimension as the coefficient of the logarithm]
With $\log\varepsilon=-2\log\Lambda_{\rm UV}$ and $\Delta=2E=-2D$, the statement $D_1=4\,(\mathrm{lc}-\{\mathcal N,-\}\ \text{units})+\dots$ is the familiar fact that the one-loop anomalous dimension is the coefficient of $\log\Lambda_{\rm UV}$ in the renormalisation of the operator, corrected for the field-strength renormalisation of its constituents (the bracket with $\mathcal N$), plus a finite gauge-fixing contribution, the composite pair kernels, and the finite counterterm of the dilatation. Here it is a theorem relating two evaluations of the same regulated graphs, and the numbers of $D_1$ come from the graphs with one regulator insertion. For translations and rotations the regulator is covariant, and their one-loop corrections are the gauge-fixing densities alone plus the finite counterterms.
\end{remark}

% ANCHOR: after_label thm:c12-local
\begin{remark}[Physics notation: locality of the gauge-fixing contribution]
In a Feynman-gauge computation the graphs with $\delta_x$ inserted on an internal ghost--gauge line split into graphs through two constituent fields and self-energy-type graphs on one field; the latter contain the composite two-point function, nonlocal through $\log p^2$ on $A_m$, $\lambda$, $\bar\lambda$ and $c$. For a BRST-invariant operator the Slavnov--Taylor identity cancels these nonlocal terms between the two kinds of graphs, the brackets $\pm\epsilon_\Xi\{\Xi_x,O\}$ of the proof, and the whole insertion becomes the sum over pairs of constituent fields of the finite parts of the graphs with one regulator insertion. For a gauge-variant operator the nonlocal terms survive in the gauge-fixing contribution and cancel in the full one-loop generator against the other graphs of Theorem \ref{thm:c12-master}.
\end{remark}

% ANCHOR: after_label thm:c12-F0
\begin{remark}[Physics notation: the scalar bilinears]
The scalar bilinears $\operatorname{tr}A_IA_J$ split into the twenty traceless ones, the $\mathbf{20}'$ of $SO(6)$, which are half-BPS, and the trace $\operatorname{tr}A_IA_I$, the Konishi operator. $D_1=0$ on the $\mathbf{20}'$ and $D_1=-48$ units $=\tfrac{3g^2N}{2\pi^2}=2\gamma_K$ on the Konishi operator, $\gamma_K=\tfrac{3g^2N}{4\pi^2}$, are the scheme-independent content. $\mathcal B_1$ is the classical supersymmetry variation of the bilinear, $\delta_\xi\operatorname{tr}A_IA_J\propto\operatorname{tr}(A_{(I}\,\xi\lambda)_{J)}$, and $Q_1=c_Q\mathcal B_1$ on the $\mathbf{20}'$ rescales it by $1+c_Q$ units; the terms $\mathcal B_2$, $\mathcal B_3$ act on the Konishi operator, with the whole connection $A=A_m+A_I$ multiplying $\xi\lambda$. The coefficients $c_Q$, $c_S$, the translation $P_1$ and the terms in $\log\tfrac32$ depend on the normalisation of the operators and of the generators; $P_1=0$ and $M_1=0$ hold at $t_P=t_M=-\tfrac13$, where translations and rotations act classically.
\end{remark}

% ANCHOR: after_label prop:c12-twist
\begin{remark}[Physics notation: the twist on the chiral bilinears]
For the supercharge of rank one with parameter $\xi_\alpha{}^i=u_\alpha z^i$, the scalars with $w_{ij}z^j=0$ are the three complex scalars $\gamma^I=A_{I4}$ for $z=e_4$, and $\operatorname{tr}\gamma^{(I}\gamma^{J)}$ are the chiral primaries of the $\mathbf{20}'$ in the holomorphic twist. Proposition \ref{prop:c12-twist} states that their one-loop twisted differential vanishes identically, in every scheme and for every normalisation of the map to the twisted letters: at one loop they stay in the cohomology of the twisted supercharge, as the half-BPS bilinears must, which is the two-letter case of [BK, (3.77)].
\end{remark}

% ANCHOR: after_label prop:c12-spectra
\begin{remark}[Physics notation: spin chains and harmonic numbers]
On single-trace operators $D_1=-8$ units times the eigenvalue of the nearest-neighbour Hamiltonian $\sum_lH_{l,l+1}$, which is $\tfrac{g^2N}{4\pi^2}$ times that eigenvalue, twice the one-loop anomalous dimension $\gamma=\tfrac{g^2N}{8\pi^2}E$. On the $\mathfrak{su}(2)$ sector of two complex scalars $Z$, $X$ the Hamiltonian is the Heisenberg chain $\sum_l(1-P_{l,l+1})$ with the permutation $P$; on $\mathfrak{su}(1|1)$ it is one minus the graded permutation; on twist two, $\operatorname{tr}(D_n^kZ\,D_n^{S-k}Z)$ with $D_n$ a light-cone covariant derivative, the eigenvalues are $4h(S)$ with the harmonic numbers $h(S)=\sum_{k=1}^S\tfrac1k$, so $\gamma=\tfrac{g^2N}{2\pi^2}h(S)$; $\operatorname{tr}(F_{11}F_{11})$ has the Konishi value. These are the one-loop spectra of [B03c].
\end{remark}
